
\chapter{Systèmes d'alimentation}
\label{chap7}


Dans ce module, nous décrirons comment les différents types de déchets sont injectés dans les installations de combustion des déchets dangereux.

\section{Injection de déchets}

\begin{figure}[H]
\centering
\includegraphics[scale=0.7]{Images/M7/M7_1.png}
\end{figure}


Le schéma ci-dessus résume les différents type de déchets et leur mode d’injection dans le four ou la post combustion.

Les solides et les boues préparés (broyés, mélangés) sont stockés dans les fosses et sont injectés dans le four par :

\begin{itemize}
   \item 	la goulotte à solides
   \item 	la pompe ou la vis à boues
\end{itemize}

\vspace{0.3cm}

Les liquides sont stockés en fonction de leurs caractéristiques dans des stockages dédiés (HPC, MPC, BPC, Chlorés…). Chaque stockages est raccordé au four par une ligne et la façade du four comprend plusieurs cannes d’injection (4 à 6) par lesquelles sont pulvérisés les déchets liquides. Certains liquides (eaux polluées) peuvent être injectés de la même façon en Post Combustion (PC).


Les conditionnés peuvent être prétraitrés (Module 4’) pour rejoindre une des filière d’injection vrac (liquide ou solide).
Certains conditionnés très réactifs ne peuvent être prétraités, ni mélangés. Ils sont alors injectés dans le four par deux techniques dites d’« injection directe » :



\begin{itemize}
\item 	injection directe solide : les déchets sont introduits directement dans le four par un sas prévu dans la goulotte à solides.
\item 	injection directe liquide : les liquides spéciaux sont pompés, ou poussés à l’azote depuis leur contenant (généralement isocontainer) vers la façade à travers des lignes dédiées.
\end{itemize}


\begin{figure}[H]
\centering
\includegraphics[width=\textwidth]{Images/M7/M7_2.png}
\end{figure}




\section{Les brûleurs}


Le brûleur permet le mélange dans les bonnes proportions du combustible (le fioul ou le gaz) et du comburant (l’air).
Il assure l’inflammation du mélange et le maintien de la combustion.
Le brûleur comporte des systèmes de sécurité qui, en cas de problèmes (absence de flamme, par exemple), doivent arrêter l’alimentation en combustible et signaler le défaut.


\subsection{Le rôle des brûleurs en incinération}

Le brûleur installé sur le four d’incinération est mis en route pour les raisons suivantes :

\begin{itemize}
\item lors des phases de démarrage du four;
\item pour assurer le soutien du four en cas de chute de la température du foyer;
\item lors des phases d’arrêt.
\end{itemize}

Au démarrage, le brûleur réchauffe le four à une température de 850°C avant l’introduction des déchets dangereux.
\vspace{0.3cm}

Le conducteur du four doit respecter des paliers de montée en température. Ces paliers sont indiqués par le fabricant du four, ils sont de l’ordre de 30 à 50°C par heure sur les gaz de combustion. Ils évitent les fortes contraintes thermiques qui risqueraient d’endommager les réfractaires et les différentes parois de la chaudière.
Dès que la température du foyer est supérieure à 850°C, le conducteur peut alors commencer à introduire les déchets dans le four. La puissance du brûleur est progressivement diminuée au profit de celle dégagée par la combustion des déchets. Le conducteur du four s’assurera que la température du foyer respecte la règle des 850°C pendant deux secondes et que le débit de vapeur reste stable.

\vspace{0.3cm}

Lors des phases de soutien, le brûleur démarre ou s’arrête automatiquement lorsque la température du foyer est inférieure ou supérieure à 850°C deux secondes.
Lors de l’arrêt complet du four, le brûleur est maintenu en marche et assure une température de 850°C jusqu’à ce qu’il n’y ait plus d’ordures en feu présentes dans le four.
\vspace{0.3cm}

Les puissances des brûleurs installés en incinération représentent 75\% de la puissance du four.



\subsection{Types de brûleurs en incinération}


On rencontre plusieurs sortes de brûleurs, que l’on peut répertorier selon :
\begin{itemize}
\item  la nature du combustible utilisé;
\item  la technologie.
 

\end{itemize}

Les brûleurs peuvent être alimentés en :
\begin{itemize}
\item  fioul domestique, ou FOD;
\item  fioul lourd (rarement utilisé);
\item  gaz naturel.

\end{itemize}


Compte tenu des températures élevées qui règnent à l’intérieur du four, les fabricants proposent différents systèmes qui permettent de protéger le brûleur :

\begin{itemize}
\item  brûleur à tête rétractable et dotée d’une guillotine (ou bouclier thermique);
\item  brûleur à tête pivotante;
\item  brûleur fixe refroidi par air.

\end{itemize}

Cette liste n'est pas exhaustive.
 
 
 
 
\subsection{Description des différents types de brûleurs}


\textbf{Technologie commune à tous les brûleurs}
\vspace{0.3cm}

Pour créer une flamme, il faut mélanger l’air au combustible et apporter de la chaleur pour provoquer l’inflammation. L’apport de chaleur est généralement provoqué par un arc électrique. On trouve aussi des torches d’allumage alimentées au propane.
 
Les brûleurs sont composés principalement de quatre circuits principaux :

\begin{itemize}
\item  le circuit d’amenée d’air de combustion;
\item  le circuit d’amenée du combustible;
\item  le circuit de mélange, c’est-à-dire l’endroit où est réalisée la combustion et où se développe et se stabilise la flamme;
\item  le circuit électrique qui assure l’inflammation du mélange, la détection de flamme, l’alimentation des moteurs de ventilateurs et des électrovannes.

\end{itemize}



\vspace{0.3cm}
\textbf{Brûleurs fonctionnant au FOD}
\vspace{0.3cm}

Description du brûleur escamotable avec écran thermique et système de pulvérisation mécanique

\begin{figure}[H]
\centering
\includegraphics[width=0.8\textwidth]{Images/M7/M7_3.png}
\end{figure}




Les éléments principaux qui constituent le brûleur sont :
\begin{itemize}
  
\item  la tête rétractable du brûleur : ce système permet d’éloigner l’ensemble des éléments de la tête de combustion du four lors de l’arrêt du brûleur, ou permet d’introduire l’ensemble des éléments dans le four lors de la mise en route. Le système est mû par un vérin alimenté en air comprimé;
\item  les manomètres qui permettent de visualiser les valeurs des pressions des différents fluides (air comprimé, air de combustion, fioul);
\item  la pompe fioul qui porte le fioul à la pression de pulvérisation (ou d’atomisation);
\item  la volute d’air qui apporte l’air de combustion au brûleur de manière homogène autour de la tête de combustion;
\item  la boîte de régulation hydraulique qui permet le réglage de la puissance du brûleur en dosant le débit de fioul et celui de l’air de combustion;
\item  le bornier de raccordement électrique;
\item  la guillotine : elle constitue un écran protégeant le brûleur des rayonnements de la flamme produite par les déchets; elle évite le phénomène de cokéfaction du fioul, c’est-à-dire la transformation du fioul en un charbon très dur sous l’effet des températures élevées. La cokéfaction se forme au niveau de la tête de combustion et empêche la combustion complète du fioul lorsque le brûleur fonctionne.

  
\end{itemize}

Les schémas suivants présentent deux situations de l’état du brûleur :
\vspace{0.3cm}

\textbf{a/ pendant l'utilisation}

\vspace{0.3cm}

Lorsque le brûleur est en marche, la tête de combustion est avancée dans le four et la guillotine de protection levée.

\begin{figure}[H]
\centering
\includegraphics[width=0.6\textwidth]{Images/M7/M7_4.png}
\end{figure}

\vspace{0.3cm}
\textbf{b/ à l’arrêt}
\vspace{0.3cm}

Lorsque le graveur est arrêté , la tête de combustion recule et le bouclier thermique s'abaisse.

\begin{figure}[H]
\centering
\includegraphics[width=0.6\textwidth]{Images/M7/M7_5.png}
\end{figure}

Le circuit combustible
Le brûleur est composé d’un circuit combustible qui comprend :

\begin{itemize}
\item  une pompe;
\item  une électrovanne;
\item  un filtre;
\item  un gicleur.

\end{itemize}

La pompe alimente le brûleur en fioul avec un débit constant et une pression élevée supérieure à la pression de pulvérisation.
\vspace{0.3cm}

La combustion du fioul est facilitée lorsque celui-ci est pulvérisé en fines gouttelettes et forme un brouillard. La pulvérisation a aussi pour rôle de faciliter le mélange du fioul avec l’air de combustion.

\vspace{0.3cm}

Les filtres placés avant la pompe et le gicleur retiennent les poussières et évitent la destruction de la pompe et l’encrassement du gicleur.
Le gicleur est équivalent à un pulvérisateur.

\begin{figure}[H]
\centering
\includegraphics[width=0.8\textwidth]{Images/M7/M7_6.png}
\end{figure}

\vspace{0.3cm}
La buse est un atomiseur.


\begin{figure}[H]
\centering
\includegraphics[width=0.5\textwidth]{Images/M7/M7_7.png}
\end{figure}


Pour que la pulvérisation se fasse correctement, la pression du fioul au niveau du gicleur doit être supérieure à 8 bars. Cette pression peut monter jusqu’à 30 bars.
Si la pression est inférieure à 7 bars, il ne se forme pas de brouillard à la sortie du gicleur, mais un filet qui ne permet pas la formation de la flamme.


\vspace{0.3cm}
\textbf{Le circuit d’air de combustion}
\vspace{0.3cm}

Il commence à l’ouïe d’aspiration et se termine au niveau du mélange air-gaz. Tout ou une partie de l’air peut servir à un prémélange air-gaz en amont de la tête du brûleur. 
Il est composé :

\begin{itemize}
\item  d’un ventilateur;
\item  d’un registre d’air;
\item  de ventelles.
\end{itemize}


Pour les grosses puissances, le ventilateur est généralement situé en dehors du brûleur, ceci pour limiter l’encombrement et le niveau de bruit. Il doit injecter le débit d’air nécessaire à la combustion du fioul et doit souffler l’air à une pression suffisante afin de vaincre les pertes de charges du circuit d’air.
\vspace{0.3cm}


\vspace{0.3cm}

Tous les combustibles (solides, liquides et gazeux) ont des caractéristiques physiques et thermiques différentes.
Trois combustibles sont utilisés par les brûleurs dans four d’incinération : le fioul domestique et le gaz, utilisés pour le brûleur principal; le propane, utilisé pour les torches d’allumage du brûleur principal.



\begin{figure}[H]
\centering
\includegraphics[width=0.7\textwidth]{Images/M7/M7_14.png}
\end{figure}


Le fioul est essentiellement composé de carbone (symbole C) et d’hydrogène (symbole H).
À ceci s’ajoute des produits non éliminés au raffinage, tel le soufre.

\vspace{0.3cm}

Le gaz naturel est constitué d’hydrocarbures, des corps composés de carbone et d’hydrogène. On trouve par exemple du méthane, composé d’un carbone lié à quatre hydrogènes (symbole CH4) qui représente 97,3\% de la composition du gaz naturel.


\textbf{Caractéristiques des combustibles}
\vspace{0.3cm}

\textbf{Masse volumique et densité}
\vspace{0.3cm}

Le fioul a une masse volumique égale à 840 kg/m3 et une densité inférieure à celle de l’eau. Ces caractéristiques sont données pour une température de 15°C.
Le gaz naturel a une masse volumique égale à 0,74 kg/m3. Il est plus léger que l’air puisque sa densité est égale à 0,57.
\vspace{0.3cm}

Ces caractéristiques sont importantes car elles permettent de prendre certaines précautions lors de l’utilisation de ces produits.
\vspace{0.3cm}

Pour le fioul, on retrouve au fond de la cuve l’eau et les boues.
Pour le gaz naturel, plus léger que l’air, il est nécessaire d’aménager des ventilations hautes dans les bâtiments fermés, ceci pour évacuer les fuites de gaz éventuelles.
\vspace{0.3cm}

Le propane, plus lourd que l’air, a tendance à se répandre au sol en cas de fuite. Il faut une extraction en bas du bâtiment.

\vspace{0.3cm}
\textbf{Viscosité}
\vspace{0.3cm}

La viscosité ne concerne que le fioul. Elle représente la résistance à l’écoulement dans les canalisations et joue un rôle important dans les opérations de pompage et de pulvérisation.
\vspace{0.3cm}

Plus le fluide est visqueux, plus il est difficile de le pulvériser et de le faire circuler dans les canalisations. La viscosité diminue avec l’élévation de température. On trouve parfois sur les lignes de gicleurs des systèmes de réchauffage de fioul qui facilitent l’écoulement et la pulvérisation du fioul. L’unité de la viscosité est le centistockes (cSt). Le fioul domestique à 20°C a une viscosité égale à 8 cSt.

\vspace{0.3cm}
\textbf{Pouvoir calorifique inférieur des combustibles}
\vspace{0.3cm}

Les quantités d’énergie calorifique (ou pouvoir calorifique) dégagées par la combustion complète d’une unité de combustible sont les suivantes :


\begin{figure}[H]
\centering
\includegraphics[ width=0.7\textwidth]{Images/M7/M7_15.png}
\end{figure}



NB : pour les gaz, l’unité de combustible est le normal mètre cube (Nm3) (voir module 5).
Puissance théorique du brûleur

\vspace{0.3cm}
À partir de ces valeurs, on peut calculer la puissance théorique dégagée par un brûleur.
Par exemple, un brûleur qui débite 1200 kg/h de fioul domestique a une puissance de :

\[
	1200 \times 11,9 = 14 280 \text{kW}
\]

\[
\text{Puissance brûleur} = \text{débit combustible} \ times \text{NCV}
\]

\textbf{Le principe de la combustion}
\vspace{0.3cm}

Il ne peut y avoir de flamme que si les trois éléments suivants sont en présence :
\begin{itemize}
\item  le combustible (le fioul ou le gaz naturel);
\item  le comburant (l’air);
\item  la chaleur.

\end{itemize}


Pour réaliser la combustion d’un mélange de combustible et d’air, il faut :
\begin{itemize}
\item  porter un point du mélange à une température suffisante ou température d’inflammation. Cette température varie de 350 à 650°C selon les combustibles. Cette température peut être atteinte par différents moyens sur le brûleur, soit par l’utilisation d’un arc électrique, soit par celle d’une torche.

\item  réaliser un mélange d’air et de combustible dans des proportions bien précises.
Pour brûler 1 Nm3 de gaz naturel dans les conditions stoechiométriques, il faut apporter 10 m3 d’air.
Pour brûler 1 kg de fioul dans les conditions stoechiométriques, il faut apporter 11,4 m3d’air.
La combustion est réalisée avec un excès d’air de l’ordre de 10%.

\end{itemize}

\subsection{La sécurité}

\vspace{0.3cm}

\textbf{Consignes communes quel que soit le type de brûleur}
\vspace{0.3cm}


Vérifiez périodiquement l’état des canalisations d’amenée du combustible, ainsi que les appareillages.
Il est nécessaire de manœuvrer régulièrement les vannes d’isolement pour éviter leur grippage. Une vanne en mauvais état doit être signalée.
Lors du démarrage du brûleur, il est conseillé de ne pas rester en face.
\vspace{0.3cm}

En cas d’arrêt volontaire des brûleurs pour des travaux :

\begin{itemize}
\item  coupez l’alimentation en électricité du brûleur;
\item  coupez la vanne générale d’alimentation de combustible;
\item  indiquez cette coupure par une pancarte.

\end{itemize}


Les extincteurs à proximité du brûleur doivent être accessibles et en bon état de marche. Ils doivent faire l’objet de contrôle annuel par un organisme.
En cas d’incendie, respectez les procédures propres à chaque site et prévenez les pompiers.


\vspace{0.3cm}
\textbf{Risques liés au gaz naturel et propane}
\vspace{0.3cm}

En quantité importante, le gaz naturel présente un risque d’asphyxie pour l’individu qui le respire. Étant inodore, le fournisseur lui ajoute un gaz odorant afin de mieux le détecter.
\vspace{0.3cm}

Dans une certaine proportion avec l’air (environ 10\%), le gaz naturel peut provoquer une explosion en présence d’une étincelle ou d’un point chaud.
En cas de fuite, il faut respecter au minimum les consignes suivantes :

\begin{itemize}
\item  à la moindre odeur de gaz, le signaler immédiatement et, si possible, isoler le circuit;
\item  aérer le local où se situe la fuite.

\end{itemize}


Il existe des appareils spéciaux pour détecter l’endroit précis de la fuite.
Il est interdit de mettre en marche des appareils électriques.
Lorsque l’on entreprend des travaux sur une ligne de gaz, il faut respecter les procédures.
\vspace{0.3cm}

Il est interdit de fumer, d’effectuer des travaux de soudage, de meulage ou perçage à proximité des zones où il y a présence de gaz.
La vérification de la pression de service permet de contrôler qu’il y a présence de gaz dans les canalisations.



 
\section{Les injecteur de déchets Vrac Liquides}

Les injecteurs de liquides, encore appelés lances ou canne d’injection font partie intégrante de l’équipement de chauffe (communément appelé « brûleur »).
Un « brûleur » à déchets comprend donc un injecteur central de combustible (fuel ou gaz ; décrit au § précédent) et plusieurs lances d’injection (4 à 6).
\vspace{0.3cm}

Chaque lance peut être raccordée à une ligne déchet, elle même reliée à une cuve.
Cette ligne d’amener du déchet peut être simple ou en boucle, c’est à dire avec un retour du déchet vers la cuve.


\begin{figure}[H]
\centering
\includegraphics[width=0.7\textwidth]{Images/M7/M7_16.png}
\end{figure}

Les injecteurs sont constitué d’un corps en acier par lequel est amené le déchet liquide et d’une double enveloppe par laquelle est amené l’air de pulvérisation (5-7 bars).
La pulvérisation se fait dans un gicleur (ou buse).
\vspace{0.3cm}

Les lignes d’injection des liquides sont équipées de débitmètres massiques qui mesurent instantanément la quantité (en kg/h) de liquides injectés dans le four et la post-combustion.
\vspace{0.3cm}

La pression d’injection du déchet ainsi que la pression de pulvérisation également suivi en continu, pour détecter les bouchages de ligne.
Une purge à l’air comprimé ou à la vapeur permet de déboucher et de nettoyer les lignes.

\begin{figure}[H]
\centering
\includegraphics[width=0.7\textwidth]{Images/M7/M7_17.png}
\end{figure}


\section{Le grappin et la goulotte d’injection des solides}

\subsection{Le pont roulant et son grappin}

\textbf{Roles}

\vspace{0.3cm}

Il sert à alimenter le four en déchets depuis la fosse jusqu’à la trémie. Il permet également de mélanger les déchets dans la fosse de manière à avoir un déchet plus homogène. Compte tenu de son utilisation intensive, jusqu’à dix mouvements par heure en plus de l’opération d’homogénéisation, c’est un équipement soumis à de nombreuses contraintes et qui peut donc connaître certaines pannes. Il est important d’en prendre le plus grand soin et de l’utiliser avec une grande attention.

\begin{figure}[H]
\centering
\includegraphics[width=0.4\textwidth]{Images/M7/M7_18.png}
{La grue au-dessus de la fosse}
\end{figure}

\textbf{Description}

\vspace{0.3cm}

Le pont roulant permet le déplacement du grappin et de la charge selon trois axes :
\begin{itemize}
\item axe de translation (sens de la longueur de la fosse);
\item axe de direction (sens de la largeur de la fosse);
\item axe de levage.
\end{itemize}

Les ponts roulants utilisés pour la manutention des déchets ménagers sont des ponts roulants bipoutres posés. Ils se déplacent en translation au-dessus de la fosse par un système de quatre galets motorisés roulant sur des rails.

\begin{figure}[H]
\centering
\includegraphics[width=\textwidth]{Images/M7/M7_19.png}
\end{figure}

Un pont roulant est composé des éléments suivants :
\begin{enumerate}
 \item une passerelle d'accès utilisée à des fins d'entretien ;
 \item la structure de la grue comprend deux poutres et quatre rouleaux; deux sont équipés de moteurs électriques et de freins intégrés ;
 \item le chariot directionnel se déplace le long des deux poutres transversales ;
 \item le grappin;
 \item plusieurs câbles d'alimentation pour moteurs et commandes.
\end{enumerate}

\begin{figure}[H]
\centering
\includegraphics[width=\textwidth]{Images/M7/M7_20.png}
\end{figure}

Le système de levage à grappin est monté sur le chariot. Le système comprend :
\begin{itemize}
\item	de quatre galets de direction; deux d’entre eux sont équipés de moteurs électriques et de freins intégrés;
\item		du moteur de levage avec ou sans variateur de vitesse, accouplé au tambour de levage à deux enroulements symétriques;
\item		d’un frein à disque et de pinces qui servent à bloquer le grappin en position d’attente lorsque la manœuvre de levage est interrompue;
\item		d’un moto-réducteur accouplé à l’enrouleur de câble d’alimentation du grappin uniquement dans le cas d’un grappin électro-hydraulique;
\item		parfois d’un système intégré de pesage des ordures contenues dans le grappin.
 
\end{itemize}


\begin{figure}[H]
\centering
\includegraphics[width=0.5\textwidth]{Images/M7/M7_21.png}
{Trolley}
\end{figure}

\textbf{Caractéristiques}

\vspace {0.3cm}

Le constructeur prescrit les principales caractéristiques de son pont. Celles-ci permettent de définir les limites de fonctionnement du pont roulant.
On trouvera inscrit :

\begin{itemize}
\item	la charge maximale ou force de levage, c’est-à-dire le poids des déchets additionné à 
celui du grappin que le pont peut soulever, indiquée en daN;
\item	la vitesse maximale des différents déplacements, à savoir, la vitesse de levage, de translation, de direction indiquée en m/mn.
\end{itemize}

Exemple de données :

Caractéristiques techniques de la grue

\begin{figure}[H]
\centering
\includegraphics[width=0.5\textwidth]{Images/M7/M7_22.png}
\end{figure}

Le fabricant indiquera la capacité ou le volume du grappin en m3

\subsection{Sécurité}

Pour éviter la rupture des câbles et du pont, il est impératif de ne pas dépasser la charge maximale admissible indiquée par le constructeur.
\vspace{0.3cm}

Il est conseillé d’utiliser le grappin en position verticale pour éviter les sorties de câble des poulies.
\vspace{0.3cm}

Il faut éviter de heurter les murs de la fosse et de provoquer des à coups.

\begin{figure}[H]
\centering
\includegraphics[width=0.5\textwidth]{Images/M7/M7_23.png}
\end{figure} Il

Le pont roulant est équipé d’un ensemble de capteurs qui limitent les dysfonctionnements et les accidents, notamment :
\begin{itemize}
\item	contrôle du mou de câble pour éviter le décrochage des charges : il stoppe automatiquement la descente du grappin lorsque celui-ci touche le sol;
\item	contrôle de chevauchement de câble pour éviter la détérioration du câble;
\item	contrôle anticollision activé lors de l’utilisation d’un deuxième pont ou lors de l’approche des murs d’enceinte de la fosse;
\item	contrôle de surcharge pour éviter la rupture du pont.
\end{itemize}

Plusieurs types de grappins existent :
\begin{itemize}
\item	les grappins à câbles acier;
\item		les grappins électro-hydrauliques (système hydraulique embarqué sur le grappin).

 \end{itemize}
Et différentes formes :
\begin{itemize}
\item		Grappin à solides (à 3 ou 4 crocs)
\item		Grappin à boues (à deux mâchoires étanches)
\end{itemize}

\begin{figure}[H]
\centering
\includegraphics[width=\textwidth]{Images/M7/M7_24.png}
\end{figure}

\vspace{0.5cm}
\textbf{Dysfonctionnements possibles}
\vspace{0.3cm}

\begin{figure}[H]
\centering
\includegraphics[width=0.7\textwidth]{Images/M7/M7_25.png}
\end{figure}


\textbf{La grue}
\vspace{0.3cm}

Certains dysfonctionnements peuvent être rencontrés au niveau des pièces en mouvement car elles sont soumises à :
\begin{itemize}
\item	une usure prématurée des galets et des rails de translation et de direction;
\item		une usure normale des câbles de levage du grappin;
\item		une usure normale des freins à disque et garnitures des pinces;
\item		une usure du tambour de levage.
\end{itemize}

Le pont roulant travaille dans un hall qui est rapidement recouvert d’une poussière très abrasive et qui provoque ainsi une usure rapide de ces pièces. Il convient de nettoyer complètement le pont une fois par semaine.
D’autres pannes peuvent survenir pour des raisons mécaniques.
\vspace{0.3cm}

Les câbles de levage seront changés dès que l’on verra l’un des brins détérioré.
Le pont peut s’arrêter si les câbles de levage se chevauchent au niveau du tambour d’enroulement. Il faut alors faire redescendre le grappin pour supprimer le chevauchement. Ceci est lié à la mauvaise conduite du grappin. Il faut éviter de déplacer le pont avant de lever le grappin, de même qu’il faut éviter de donner trop de mou au câble (risque de décrochement du grappin).
\vspace{0.3cm}

Le chariot peut aussi s’arrêter intempestivement alors qu’il se trouve au milieu de la fosse : ceci est dû à un mauvais fonctionnement des contacts fin de course. Il faut alors vérifier sur le pont la position des « étoiles », ou le contact de fin de course.


\begin{figure}[H]
\centering
\includegraphics[width=0.5\textwidth]{Images/M7/M7_26.png}
\end{figure}


En dehors des vérifications quotidiennes, le pontier contrôle de manière visuelle l’état 
général du pont une fois par semaine :

\begin{itemize}
\item	état du moteur ;		
\item état du câble de levage ;
\item	état du câble d'alimentation/de commande ;	
\item état des freins (plaquettes et disques);
\item état des	rouleaux et des rails ;	
\item cellules de détection anticollision.
\end{itemize}

\vspace{0.3cm}
\textbf{Le grappin}

\vspace{0.3cm}

Il peut arriver également que l’ouverture ou la fermeture du grappin ne fonctionne pas correctement à cause de déchets coincés entre les dents (ou les vérins de commande dans le cas du système électro-hydraulique). À chaque prise de quart, il faut donc vérifier les endroits « sensibles » du grappin et procéder si nécessaire au nettoyage de celui-ci.
\vspace{0.3cm}

On procédera au graissage des axes de dents au moins une fois par 24 heures, en général de nuit lorsque l’activité est moins importante.
\vspace{0.3cm}

On vérifiera visuellement qu’il n’y a pas de fuite d’huile au niveau des vérins, dans le cas de grappin électro-hydraulique. Dans ce dernier cas, on contrôlera également :

\begin{itemize}
\item	le niveau d’huile;
\item		les deux pressions de service d’huile dans le circuit hydraulique du grappin avec et sans effort de serrage, à l’aide d’un manomètre monté au refoulement de la pompe hydraulique, la pompe étant en fonctionnement.

\end{itemize}
Ces pressions sont indiquées par le constructeur (par exemple 6 et 180 bars).


\vspace{0.3cm}
\newpage
\textbf{Procédures de sécurité}

\vspace{0.3cm}

PRISE DE POSTE
\vspace{0.3cm}

Prendre connaissance des indications du pontier précédent
Vérifier que personne n’intervient sur le pont roulant
S’assurer que les accès aux quais sont verrouillés
Vérifier le bon état des câbles
Faire à vide des manœuvres de levage, de direction et de translation
Vérifier le bon fonctionnement des équipements (freins, contacts de fin de course…)

\vspace{0.3cm}
CONDUITE
\vspace{0.3cm}

Ne permettre la conduite qu’à des personnes qualifiées
Ne pas balancer le grappin pour saisir ou déposer une charge
Ne pas effectuer de traction oblique
Éviter de laisser la charge suspendue longtemps en attente dans le grappin
Interdire l’accès dans les zones où les chutes peuvent se produire
Interdire l’accès dans la zone de travail du pont
Monter le grappin au maximum avant d’entamer un mouvement de translation

\vspace{0.3cm}
FIN DE POSTE
\vspace{0.3cm}

Laisser le grappin hors d’atteinte et verrouiller le pont

\vspace{0.3cm}
ENTRETIEN COURANT
\vspace{0.3cm}

Amener le pont roulant au point prévu pour l’entretien
Consigner le pont, ouvrir l’interrupteur général, le verrouiller et interdire sa fermeture, 
le signaler avec une clé
Utiliser un équipement individuel contre les chutes (harnais, filet…)

\subsection{La goulotte des solides}

\begin{figure}[ H]
\centering
\includegraphics[width=\textwidth]{Images/M7/M7_27.png}
\end{figure}

La goulotte solides assure l’enfournement de deux types de déchets :
les déchets solides vrac depuis les fosses de stockage par l’intermédiaire du grappin
les déchets conditionnés réactifs (filière directe)
\vspace{0.3cm}

Le mouvement du grappin peut être entièrement automatisé. Le conducteur de four choisit alors par l’intermédiaire d’un programme informatique : la zone de prise dans les fosses, l’ouverture du grappin, la fréquence des mouvements…
Le grappin peut-être contrôlé en semi-automatique : le conducteur reprend la maîtrise du système au moment de l’injection dans la goulotte. C’est lui qui décide quand, et combien de déchet seront enfournés, en fonction de ses paramètres de conduite (températures, taux d’oxygène…).
\vspace{0.3cm}

L’alimentation peut également être réalisée par un tapis à partir d’un tampon. Le grappin alimente une trémie tampon qui se vide sur un tapis à écaille. Le tapis alimente une trémie sur pesons qui bascule une charge dans le four. Le conducteur de four contrôle la charge maximale par mouvement et la fréquence d’injection d’une charge.

\begin{figure}[H]
\centering
\includegraphics[width=\textwidth]{Images/M7/M7_28.png}
\end{figure}

Une fois les déchets déposés dans la goulotte, un automate gère l’ouverture et la fermeture des registres ainsi que le mouvement du vérin poussoir d’injection.

\begin{figure}[H]
\centering
\includegraphics[width=\textwidth]{Images/M7/M7_29.png}
\end{figure}



\section{Les pompes à boues}

Les pompes à boues peuvent être utilisées pour injecter une charge homogène de déchets pâteux à l'avant de l'incinérateur. Le degré de pâte peut varier. Les pompes à boues sont essentiellement des pompes à béton adaptées à cet usage spécifique.
Deux types de pompes à boues sont utilisées pour injecter les déchets dangereux :
\begin{itemize}
 \item Une pompe à piston simple : déchets pâteux épais
 \item Une pompe à double piston : déchets pâteux liquides
\end{itemize}

Les déchets alimentant ce type d'équipements sont stockés dans des fosses dédiées. Ces déchets sont soigneusement préparés par mélange avec le grappin ou un chargeur mécanique. L'opérateur préparant les déchets doit notamment veiller à ce qu'aucune pièce métallique, trop grosse pour les pompes, ne pénètre dans le mélange.
\vspace{0.3cm}

Ce matériel d'injection offre l'avantage d'injecter les déchets à la fois en douceur et régulièrement. Cela contribue à créer des conditions de combustion satisfaisantes.


\begin{figure}[H]
\centering
\includegraphics[width=\textwidth]{Images/M7/M7_30.png}
\end{figure}

\section{Les lignes directes}

Certains déchets liquides sont particulièrement réactifs et ne peuvent être mélanger sans danger avec d’autres liquides. Certains déchets très toxiques ne doivent en aucun cas être manipuler à l’air libre.
Pour ces différents types de déchets, il existe des lignes directes, c’est à dire des tuyauteries, des pompes et des lances d’injection dédiées.

\begin{figure}[H]
\centering
\includegraphics[width=0.4\textwidth]{Images/M7/M7_31.png}
\end{figure}

Généralement, les déchets concernés sont livrés en isocontainer : c’est une citerne particulière de 10 à 15 m3, équipée d’une double enveloppe et inertée à l’azote.
Dans le cas de déchet nécessitant un préchauffage avant injection, la vapeur est branchée sur l’isocontainer et il est raccordé à une ligne tracée, c’est à dire chauffé extérieurement.
L’injection se fait généralement en appliquant une pression d’azote dans l’isocontainer : le déchet est alors poussé vers la lance de pulvérisation.



\section{Injection directe} 


Les petits conditionnés réactifs, qui ne peuvent être vidés ou broyés, sont injecté directement dans le four. Il peut également s’agir de destruction sous contrôle d’huissier, de produits de Recherche et Développement (confidentialité), de drogues…
\vspace{0.3cm}

Les petits conditionnés sont tout d’abord placés dans des cartons, en fonction de leur composition chimique, stabilisés avec de la vermiculite (sable) et envoyés vers la plate-forme de « Filière Directe ».
\vspace{0.3cm}

Ils intègrent alors une chaîne de transporteurs (table à rouleaux ou tapis) qui les amène jusqu’à un sas dans la goulotte à solide. Ils sont pesés juste avant l’injection à travers ce sas, dans l’intervalle entre les deux registres de la goulotte.
\vspace{0.3cm}

Un programme informatique gère les cadence d’enfournement. Le conducteur de four sera très attentif lors de l’injection de ces déchets qui peuvent apporter des perturbation notable dans la combustion.

\begin{figure}[H]
\centering
\includegraphics[width=0.8\textwidth]{Images/M7/M7_32.png}
\end{figure}



